\scanpage{4}

\[
\begin{array}{ccc}
A_1&A_2&A_3\\
1\ u_1&1\ u_2&1\ u_3
\end{array}
\qquad u_i^2+c_i=0
\]

extensions quadratiques, involutions canoniques $\sigma_1,\sigma_2,\sigma_3$ ($\sigma_i(u_i)=-u_i$) se prolongent à $A$ de façon canonique par $\sigma_1\otimes\mathrm{id}\otimes\mathrm{id}$ etc.

\[
\begin{aligned}
&A_1\otimes A_2\otimes A_3\overset{\unsure{\mathrm{déf}|\mathrm{dfn}|\mathrm{df}}}=A\\
&1\quad u_1\quad u_2\quad u_3\quad u_2u_3\quad u_3u_1\quad u_1u_2\quad u_1u_2u_3
\end{aligned}
\]

% margin note written sideways
\margin{correspondent aux sept hyperplans de $\mathbb F_2^3$, i.e. par \unsure{pol.} \unsure{orth.} aux sept droites de $\mathbb F_2^3$, i.e. aux \deleted{sept éléments non nuls}}

\diagram{$\begin{array}{llll}
k&1&\sigma_1,\sigma_2,\sigma_3,\sigma_2\sigma_3,\sigma_3\sigma_1,\sigma_1\sigma_2,\sigma_1\sigma_2\sigma_3=\tau&\\
A_1&1,u_1&\sigma_2,\sigma_3,\sigma_2\sigma_3&\\
A_2&1,u_2&\sigma_3,\sigma_1,\sigma_3\sigma_1&\\
A_3&1,u_3&\sigma_1,\sigma_2,\sigma_1\sigma_2&\\
B_1&1,u_2u_3&\sigma_1,\sigma_2\sigma_3,\sigma_1\sigma_2\sigma_3=\tau&B_1\simeq A_2\wedge A_3\\
B_2&1,u_3u_1&\sigma_2,\sigma_3\sigma_1,\sigma_2\sigma_3\sigma_{\unsure{1|3}}=\tau&B_2\simeq A_3\wedge A_1\\
B_3&1,u_1u_2&\sigma_3,\sigma_1\sigma_2,\sigma_3\sigma_1\sigma_2=\tau&B_3\simeq A_1\wedge A_2\\
B_0&1,u_1u_2u_3&\sigma_2\sigma_3,\sigma_3\sigma_1,\sigma_1\sigma_2&B_0\simeq A_1\wedge A_2\wedge A_3
\end{array}$}

\diagram{$\left.\begin{array}{lll}
A'_1=(A_{\unsure{2|3}},A_3,B_1)&1,u_2,u_3,u_2u_3&\sigma_1\\
A'_2=(A_3,A_1,B_2)&1,u_3,u_1,u_3u_1&\sigma_2\\
A'_3=(A_1,A_2,B_3)&1,u_1,u_2,u_1u_2&\sigma_3\\
B'_1=(A_1,B_1,B_0)&1,u_1,u_2u_3,u_1u_2u_3=u&\sigma_2\sigma_3\\
B'_2=(A_2,B_2,B_0)&1,u_2,u_3u_1,u_2u_3u_1=u&\sigma_3\sigma_1\\
B'_3=(A_3,B_3,B_{0})&1,u_3,u_1u_2,u_3u_1u_2=u&\sigma_1\sigma_2\\
B'_0=(B_1,B_2,B_3)&1,u_2u_3,u_3u_1,u_1u_2&\sigma_1\sigma_2\sigma_3
\end{array}\right\}$}

correspondent aux sept droites vect. de $(\mathbb F_2^3)^{\deleted{\unsure{\ast}}}$ i.e. aux \deleted{sept} sept éléments non nuls

\[
\begin{aligned}
&A_i\otimes A'_i\xrightarrow{\sim} A\\
&[B_i\subset B'_i]\\
&B_0\otimes B'_0\xrightarrow{\sim} A\quad\text{si les $c_i$ inv. i.e. les $A_i$ étales}
\end{aligned}
\]

\[
\begin{aligned}
A_1&\subset A'_2,A'_3,B'_1\\
A_2&\subset A'_3,A'_1,B'_2\\
A_3&\subset A'_1,A'_2,B'_3\\
B_1&\subset A'_1,B'_1,B'_0\\
B_2&\subset A'_{\unsure{1|2}},B'_{\unsure{1|2}},B'_0\\
B_3&\subset A'_3,B'_3,B'_0\\
B_0&\subset \underset{B_0(A_1)}{\underset{\Vert}{B'_1}},\underset{B_0(A_2)}{\underset{\Vert}{B'_2}},\underset{B_0(A_3)}{\underset{\Vert}{B'_3}}
\end{aligned}
\]

il y a $7+7=14$ situations ``\unsure{bi}quadratiques'' typiques (quand les $A_i$ sont étales\ldots)
% archive sheet number: 3
